\section*{Part III: Teaching Is a Medium-Specific Craft}

\subsection*{20. Start with the Function, Not the Platform}

Founders frequently begin by asking, ``Should the course be online?'' or ``Should every lecture be live?'' That is the wrong 1st question. A medium is an arrangement of affordances and constraints. A particular learning activity has a cognitive and social function. The design problem is to match them.

Reading supports precision, revisiting, and inspecting formal arguments. A recorded demonstration supports pausing and replaying a process. Live online sessions support remote conversation and real-time review of shared digital work. A physical laboratory supports equipment, embodied procedures, and sensory constraints. Independent projects support integration and autonomy. Peer discussion supports articulation and confrontation with alternative reasoning. A well-trained human mentor can diagnose mistakes that a constrained auto-grader cannot; an auto-grader can identify routine test failures more cheaply and quickly than a human.

The question is \textbf{which medium delivers the necessary function at a tolerable total cost}, including preparation, accessibility, learner time, and review capacity. Mixed delivery is not automatically superior because it mixes formats. It is superior only if the combination makes the learning mechanism stronger or the experience more feasible.

\subsection*{21. The Classroom: Use the Room for Shared Thinking}

A physical classroom provides proximity, spontaneous interaction, rich nonverbal cues, shared attention, and the possibility of working with objects or equipment. None of these advantages arises merely because learners sit in the same place. A lecturer speaking for 2 hours while 30 students transcribe slides uses the room largely as an expensive broadcasting channel.

Design an in-person class around what physical co-presence makes possible: simultaneous individual attempts, visible comparisons, short debates about competing solutions, rapid diagnosis of common errors, collaborative physical or digital experiments, and a final independent check. A room should permit learners to see the relevant work and participate without requiring extroversion. Seating, acoustics, power outlets, accessible entry, readable projection, and ventilation are instructional design decisions, not facilities trivia.

Consider a 90-minute session on evaluation validity. The 1st 10 minutes retrieve prerequisites and ask everyone to judge an apparently excellent model. The teacher then introduces the data-generating problem with a diagram. Students independently inspect 1 suspect pipeline before discussing it in pairs. The instructor samples explanations and models a diagnostic method. Groups compare 3 cases that look similar but require different splits. Everyone submits a final unseen decision. The exact timings are flexible; the enduring design is \textbf{independent judgment $\to$ explanation $\to$ feedback $\to$ varied reapplication}.

\subsection*{22. Discussion Should Have a Cognitive Contract}

``Discuss with your neighbor'' can be a very weak instruction. Students may chat, exchange guesses, or let the most confident person decide. Specify the problem, the artifact, the disagreement that matters, and the individual accountability after collaboration. For example: ``Each of you must 1st choose a validation plan. Compare why you chose it. As a pair, identify 1 assumption that could reverse the decision. Both of you will then answer a new case independently.''

For group projects, establish roles that rotate: analyst, skeptic, implementer, reviewer, and presenter. Ensure that no person can routinely pass through the experience without attempting the target reasoning. Peer learning works when the peers actively make 1 another's thinking more explicit. It fails when the institution mistakes physical proximity for intellectual collaboration.

\subsection*{23. The Teacher Needs a Branching Plan, Not Only a Script}

In-person instruction benefits from responsiveness. A rigid monologue assumes the class will misunderstand precisely when predicted. An effective teacher prepares several branches. If most participants fail a prerequisite check, show a smaller worked example. If a minority fail, provide targeted support while the others attempt an extension. If most learners already understand the method, proceed to a harder counterexample. If the planned activity fails because of technology, use a paper or board version that preserves the essential cognitive work.

Write the lesson as a \emph{score of learner actions and teacher observations}, not merely the words to be spoken. Which responses would trigger a change? What errors are serious enough to interrupt the lesson? Which misconception can be deferred safely? The skill of teaching includes deciding \textbf{what not to explain right now}.

\subsection*{24. Live Online Instruction: Eliminate Invisible Attendance}

A learner can be connected to a meeting without being mentally present. Online teaching makes it easier to hide uncertainty, multitask, lose context during screen transitions, or become a silent spectator. The response is not to require cameras on as a crude proxy for learning. It is to create \textbf{small, frequent, meaningful traces of work}.

Before a live session, supply a small prerequisite task. Open with a private response rather than roll call. Share a stable document, notebook, drawing surface, or repository. Alternate concise explanation with individual problem-solving. Use breakout rooms only when they have clear outputs, appropriate sizes, and a reason that collaboration is better than individual work. Allow quiet written participation. Provide a low-bandwidth fallback and an accessible record of essential instructions.

An excellent online class can be intensely interactive without constant spectacle. For a coding workshop, learners might submit predictions about code behavior, run experiments in their own environment, compare error traces, and present 1 design trade-off. The teacher can make reasoning visible with screen annotation and selected artifacts. The live medium earns its place when it allows \textbf{immediate contingent teaching}: the instructor does something different because of what learners actually attempted.

\subsection*{25. Plan the Online Run-of-Show as a Production System}

A reliable online session needs a working plan for technology and people. Specify which screen will be visible, the link to the exact artifact, the teacher's next transition, how answers are collected, who watches the chat if a 2nd person is available, how long an activity reasonably takes, and what happens when a platform becomes unavailable.

For an initial founder-run class, many tools can be simple: a video meeting, a shared repository, 1 form, and a text-based issue log. Avoid purchasing a complex learning-management platform merely to automate an unproven process. But do rehearse transitions. 10 avoidable minutes lost to missing files can destroy a carefully balanced 60-minute session.

A teacher's after-class record should capture not only attendance and satisfaction but the actual response distribution, recurring errors, overdue feedback, and revisions required for the next session. The purpose of a run-of-show is to make adaptive instruction dependable, not theatrical.
